\documentclass[11pt,a4paper]{article}
\usepackage{amsmath,amssymb,mathtools}
\usepackage{geometry}
\usepackage{hyperref}
\usepackage{booktabs}
\usepackage{array}
\usepackage{xcolor}
\usepackage{physics}
\geometry{margin=2.5cm}

\hypersetup{
  colorlinks=true,
  linkcolor=blue,
  urlcolor=blue,
  citecolor=blue,
}

\newcommand{\PM}{\text{PM}}
\newcommand{\GR}{\text{GR}}
\newcommand{\Msun}{M_\odot}
\newcommand{\rhonuc}{\rho_\text{nuc}}
\newcommand{\rhocrit}{\rho_\text{crit}}
\newcommand{\pcrit}{P_\text{crit}}
\newcommand{\half}{\tfrac{1}{2}}
\newcommand{\muG}{\mu_G}

\title{PM Compact-Star Structure:\\Derivation, Equations, and Observational Predictions}
\author{Pushing-Medium Framework}
\date{\today}

\begin{document}
\maketitle
\tableofcontents
\newpage

%==========================================================================
\section{Overview}
%==========================================================================

This document derives the structure equations for a compact star within the
Pushing-Medium (PM) framework.  The PM scalar field $\varphi$ plays the role
of the gravitational potential, and the deformation energy of the medium
provides a self-consistent equation of state (EOS).

The key results are:
\begin{itemize}
  \item The corrected massive-particle force law $\mathbf{a} = +\tfrac{c^2}{2}\nabla\varphi$,
    which gives \emph{exactly} Newtonian gravitational acceleration at all
    compactnesses via the PM Poisson equation and Gauss's theorem.
  \item A self-consistent PM EOS $P = \tfrac{c^2}{2}(\rho - \rhonuc)$, derived
    solely from PM hydrostatics without borrowing GR results.
  \item PM stellar structure equations that are \emph{exact within PM} — not
    a weak-field approximation.  The PM and GR structure equations differ
    because they come from different physical frameworks, not because one is
    a limit of the other.
  \item A firm density ceiling $\rho_c \le \rhocrit = e\,\rhonuc$ from the
    PM deformation-energy stability criterion.
  \item A maximum-mass prediction that can be tested against current NICER
    and gravitational-wave observations.
\end{itemize}

%==========================================================================
\section{PM Sign Convention and Scalar Field}
%==========================================================================

\subsection{PM effective metric}

The PM framework models gravity as a compressed elastic medium with
refractive index $n(\mathbf{x}) = e^{\varphi}$.  The optical metric
(wave propagation metric) seen by all particles is (in isotropic
Cartesian coordinates)

\begin{equation}
  ds^2 = -c^2(1 + 2\varphi)\,dt^2 + (1 - 2\varphi)\,d\ell^2,
  \label{eq:metric}
\end{equation}

where $d\ell^2 = dx^2 + dy^2 + dz^2$.  Outside a point mass $M$ at the
origin, the field is

\begin{equation}
  \varphi(r) = \frac{\muG M}{r}, \qquad \muG \equiv \frac{2G}{c^2}
             = 1.4861 \times 10^{-27}\;\text{m}^2\,\text{kg}^{-1}.
  \label{eq:phi_exterior}
\end{equation}

This gives $\varphi > 0$ near the mass (converging medium, $n > 1$), consistent
with PM lensing predictions.

\subsection{Gradient direction near a mass}

For a mass at the origin and a test point at $\mathbf{r} = r\hat{\mathbf{r}}$:

\begin{equation}
  \nabla\varphi = -\frac{\muG M}{r^2}\hat{\mathbf{r}}
               = -\frac{2GM}{c^2 r^2}\hat{\mathbf{r}}.
  \label{eq:grad_phi}
\end{equation}

The gradient points \emph{inward} (toward the mass), i.e.\ in the $-\hat{r}$
direction.

%==========================================================================
\section{Massive-Particle Force Law}
%==========================================================================

\subsection{Light vs.\ massive particles: the factor of 2}

The well-known PM result for the light-deflection angle by a point mass $M$
(impact parameter $b$) is

\begin{equation}
  \alpha_\text{light} = \frac{4GM}{c^2 b},
  \label{eq:deflection}
\end{equation}

which requires $\varphi = 2GM/(c^2 r)$ — twice the Newtonian potential
$\Phi_N = -GM/r$ (up to sign).  This factor-of-2 arises because light is
sensitive to \emph{both} metric components: $g_{tt}$ (time dilation) and
$g_{rr}$ (spatial stretching) contribute equally.

Massive particles in the non-relativistic limit couple only to $g_{tt}$ (time
dilation alone sets the geodesic acceleration to leading order in $v/c$).
This halves the effective coupling:

\begin{equation}
  \boxed{\mathbf{a}_\PM = +\tfrac{c^2}{2}\,\nabla\varphi.}
  \label{eq:force_law}
\end{equation}

\subsection{Verification}

Inserting Eq.~(\ref{eq:grad_phi}) into Eq.~(\ref{eq:force_law}):

\begin{equation}
  \mathbf{a}_\PM = +\tfrac{c^2}{2}\left(-\frac{2GM}{c^2 r^2}\hat{\mathbf{r}}\right)
                = -\frac{GM}{r^2}\hat{\mathbf{r}}.
  \label{eq:newton}
\end{equation}

This is \emph{exactly} the Newtonian gravitational acceleration, attracting the
test particle toward the mass. $\checkmark$

\subsection{Common pitfall in PM implementations}

An incorrect sign $\mathbf{a} = -c^2\nabla\varphi$ was historically used in
some PM implementations.  With $\nabla\varphi$ pointing inward, this gives
\emph{outward} (repulsive) acceleration with magnitude $2GM/r^2$ — twice
Newtonian and in the wrong direction.  Eq.~(\ref{eq:force_law}) is the
physically correct form.

%==========================================================================
\section{PM Hydrostatic Equilibrium}
%==========================================================================

\subsection{Static PM fluid}

Consider a spherically symmetric distribution of PM medium in static
equilibrium.  The PM momentum equation in the static limit ($\mathbf{u}=0$)
for a fluid element reduces to a balance between the pressure gradient and
the PM body force:

\begin{equation}
  \nabla P = \rho(\varphi)\,\mathbf{a}_\PM = \rho(\varphi)\cdot\tfrac{c^2}{2}\,\nabla\varphi.
  \label{eq:hydrostatic}
\end{equation}

In the radial direction this is

\begin{equation}
  \frac{dP}{dr} = \rho(r)\cdot\frac{c^2}{2}\frac{d\varphi}{dr}.
  \label{eq:hydrostatic_1d}
\end{equation}

With $d\varphi/dr = -2GM_\text{enc}(r)/(c^2 r^2)$ this becomes exactly the
Newtonian hydrostatic equation:

\begin{equation}
  \frac{dP}{dr} = -\frac{G\,m(r)\,\rho(r)}{r^2}.
  \label{eq:dPdr}
\end{equation}

\subsection{Comparison with GR (TOV)}

The general-relativistic Tolman--Oppenheimer--Volkoff equation is

\begin{equation}
  \frac{dP}{dr} = -\frac{G(m + 4\pi r^3 P/c^2)(\rho + P/c^2)}{r^2(1 - 2Gm/c^2r)}.
\end{equation}

Compared to Eq.~(\ref{eq:dPdr}), the TOV equation has three extra terms:
\begin{itemize}
  \item $+4\pi r^3 P/c^2$ in the mass term (pressure as active gravitational source
    in GR),
  \item $+P/c^2$ in the density term (pressure as inertial mass in GR),
  \item $1/(1-2Gm/c^2r)$ in the denominator (spacetime curvature factor in GR).
\end{itemize}
These terms are absent in the PM structure equation \emph{because PM does not
use curved spacetime}.  Spacetime is flat in PM; the optical metric
$n = e^\varphi$ describes wave propagation in the medium, not spacetime
curvature.  Consequently there is no Schwarzschild radius, no event horizon,
and no GR compactness limit $Gm/(c^2r) < 1/2$ in PM.

The two frameworks produce different structure equations because they are
based on different physical assumptions, not because one is an approximation
to the other.  The key PM prediction distinguishing it from all GR models is
the hard density ceiling $\rhocrit$, which sets an absolute upper bound on
the central density of all stable PM compact stars.

%==========================================================================
\section{PM Self-Consistent Equation of State}
%==========================================================================

\subsection{Density–field relation}

In PM, the mass density of the compressed medium is

\begin{equation}
  \rho(\varphi) = \rhonuc\,e^{\varphi},
  \label{eq:rho_phi}
\end{equation}

where $\rhonuc \approx 2.3\times10^{17}\;\text{kg/m}^3$ is the nuclear saturation
density (medium density in the undeformed state at $\varphi = 0$).

\subsection{EOS derivation}

Combining the hydrostatic condition Eq.~(\ref{eq:hydrostatic}) with
Eq.~(\ref{eq:rho_phi}):

\begin{equation}
  dP = \rho(\varphi)\cdot\tfrac{c^2}{2}\,d\varphi
     = \rhonuc e^{\varphi}\cdot\tfrac{c^2}{2}\,d\varphi.
\end{equation}

Integrating from vacuum ($\varphi = 0$, $P = 0$) to field $\varphi$:

\begin{equation}
  P(\varphi) = \tfrac{c^2}{2}\rhonuc(e^{\varphi} - 1).
\end{equation}

Substituting $\rho = \rhonuc e^{\varphi}$:

\begin{equation}
  \boxed{P = \frac{c^2}{2}(\rho - \rhonuc).}
  \label{eq:eos}
\end{equation}

\subsection{EOS properties}

\begin{center}
\begin{tabular}{lll}
  \toprule
  Property & Value & Note \\
  \midrule
  Surface pressure & $P(\rhonuc) = 0$ & Correct surface condition \\
  Sound speed & $c_s = c/\!\sqrt{2}$ & Causal ($c_s < c$) \\
  Stiffness & $dP/d\rho = c^2/2$ & Linear (``$\Gamma = 2$''-type) \\
  \bottomrule
\end{tabular}
\end{center}

A linear EOS $P = K(\rho - \rho_0)$ with $K = c^2/2$ is comparable in
stiffness to nuclear physics polytropes SLy and APR with $\Gamma \approx
2$--$3$ above saturation density.

%==========================================================================
\section{PM Stellar Structure System}
%==========================================================================

\subsection{Structure equations}

The full PM stellar structure system in spherical symmetry is:

\begin{align}
  \frac{dm}{dr} &= 4\pi r^2\,\rho(r),  \label{eq:dmdr} \\[6pt]
  \frac{dP}{dr} &= -\frac{G\,m(r)\,\rho(r)}{r^2},  \label{eq:dPdr2}
\end{align}

with $\rho = \rhonuc + 2P/c^2$ from the PM EOS.

\subsection{Boundary conditions}

\begin{align}
  m(0) &= 0,  \\
  P(0) &= P_c = \tfrac{c^2}{2}(\rho_c - \rhonuc), \quad \rho_c \le \rhocrit, \\
  P(R) &= 0,  \quad M_\star = m(R).
\end{align}

\subsection{PM critical density constraint}

The PM critical density arises from the deformation-energy stability criterion
at $\varphi_\text{crit} = 1$:

\begin{equation}
  \rhocrit = \rhonuc\,e^1 = e\,\rhonuc \approx 6.25\times10^{17}\;\text{kg/m}^3.
  \label{eq:rho_crit}
\end{equation}

No stable PM compact star can have a central density exceeding $\rhocrit$.
This is an absolute PM prediction.  The corresponding central pressure limit is

\begin{equation}
  \pcrit = \half c^2(\rhocrit - \rhonuc) = \half c^2\,\rhonuc(e-1)
         \approx 1.87\times10^{34}\;\text{Pa}.
\end{equation}

%==========================================================================
\section{Observational Predictions and Kill Criteria}
%==========================================================================

\subsection{PM M–R curve}

Sweeping the central density $\rho_c$ from $\rhonuc$ to $\rhocrit$ and
integrating the structure system Eqs.~(\ref{eq:dmdr})--(\ref{eq:dPdr2})
generates the PM mass–radius curve.

The maximum-mass PM star occurs at $\rho_c = \rhocrit$.  The predicted
$M_\text{max}$ and $R$ at $M_\text{max}$ depend on the stiffness of the PM EOS
and the PM structure equations (exact within PM).

\subsection{Comparison to GR / TOV}

The PM and GR structure equations differ because they come from different
physical frameworks (flat medium vs.\ curved spacetime).  The observable
consequences are:
\begin{itemize}
  \item PM gives \emph{larger radii} than GR for the same central density,
    because the GR metric factor $1/(1-2Gm/c^2r)$ (absent in PM) compresses
    the pressure gradient and shrinks the star.
  \item PM imposes a \emph{hard density ceiling} $\rhocrit = e\,\rhonuc$ with
    no GR analogue.  This limits the peak of the PM M\u2013R curve independently
    of EOS stiffness.
\end{itemize}

\subsection{Key observational tests}

\begin{center}
\begin{tabular}{p{4.5cm} p{4.5cm} p{4.5cm}}
  \toprule
  Observation & Constraint & PM implication \\
  \midrule
  PSR J0952-0607 & $M = 2.35 \pm 0.17\,\Msun$ & PM must predict $M_\text{max} \ge 2.18\,\Msun$ \\[4pt]
  PSR J0740+6620 (NICER) & $M = 2.08\,\Msun$, $R \approx 12.4\,$km & Check PM M–R passes through box \\[4pt]
  PSR J0030+0451 (NICER) & $M \approx 1.34\,\Msun$, $R \approx 12.7\,$km & Check PM M–R passes through box \\[4pt]
  GW170817 & $R_{1.4} \approx 11.9\,$km & PM must predict $R_{1.4}$ in this range \\
  \bottomrule
\end{tabular}
\end{center}

\subsection{PM falsification criterion}

\begin{description}
  \item[Kills PM if:] The PM M–R curve does not pass through the NICER
    observational boxes, OR $M_\text{max} < 2.18\,\Msun$ (PM cannot support
    the heaviest known pulsar).

  \item[Supports PM if:] The PM M–R curve passes through both NICER boxes and
    predicts $R_{1.4}$ within the GW170817 window, with
    $M_\text{max} \ge 2.18\,\Msun$.
\end{description}

The PM structure equations integrated here are exact within the PM framework.
No GR metric corrections are missing or approximated.  The M\u2013R curve
computed by \texttt{solve\_pm\_star} is the genuine PM prediction.

%==========================================================================
\section{Implementation Notes}
%==========================================================================

The stellar structure equations are implemented in
\texttt{src/pushing\_medium/stellar\_structure.py}:

\begin{itemize}
  \item \texttt{pm\_eos\_pressure(rho)} — Eq.~(\ref{eq:eos})
  \item \texttt{pm\_eos\_density(P)} — inverse EOS
  \item \texttt{pm\_eos\_sound\_speed()} — $c/\sqrt{2}$
  \item \texttt{solve\_pm\_star(rho\_central)} — ODE integrator using
    \texttt{scipy.integrate.solve\_ivp} (DOP853 method) with an event
    function to stop at $P = 0$
  \item \texttt{compute\_mr\_curve(n\_points)} — sweeps $\rho_c \in [\rhonuc,\rhocrit]$
\end{itemize}

The M--R curve script \texttt{scripts/pm\_mass\_radius\_curve.py} produces
a plot with NICER observation boxes and GW170817 constraints overlaid.

\appendix
%==========================================================================
\section{Useful PM Constants}
%==========================================================================

\begin{center}
\begin{tabular}{lll}
  \toprule
  Symbol & Value & Meaning \\
  \midrule
  $c$ & $2.998 \times 10^8\;\text{m/s}$ & Speed of light \\
  $G$ & $6.674 \times 10^{-11}\;\text{N\,m}^2\text{kg}^{-2}$ & Newton constant \\
  $\muG = 2G/c^2$ & $1.486 \times 10^{-27}\;\text{m}^2/\text{kg}$ & PM coupling \\
  $\rhonuc$ & $2.3 \times 10^{17}\;\text{kg/m}^3$ & Nuclear saturation density \\
  $\rhocrit = e\,\rhonuc$ & $6.25 \times 10^{17}\;\text{kg/m}^3$ & PM critical density \\
  $\pcrit$ & $1.87 \times 10^{34}\;\text{Pa}$ & PM central pressure limit \\
  $\varphi_\text{crit}$ & $1$ & Critical PM scalar field \\
  $c_s$ & $c/\sqrt{2} \approx 2.12 \times 10^8\;\text{m/s}$ & PM sound speed \\
  \bottomrule
\end{tabular}
\end{center}

\end{document}
